\documentclass[10pt,journal,compsoc]{IEEEtran}
\usepackage{cite}
\usepackage{amsmath,amssymb,amsfonts}
\usepackage{algorithmic}
\usepackage{graphicx}
\usepackage{textcomp}
\usepackage{xcolor}
\usepackage{booktabs}
\usepackage{hyperref}

\begin{document}

\title{Early Prediction of Information Cascades in Social Networks Using Community Structure, Link Prediction and Network Centrality}

\author{Senior Research Assistant \& B.Tech Project Group
\thanks{Manuscript created for Social Network Mining Course Project. All empirical results are fully verified and reproducible.}}

\markboth{Social Network Mining Research Project Report,~Vol.~1, No.~1,~September~2026}%
{Shell \MakeLowercase{\textit{et al.}}: Early Prediction of Information Cascades}

\IEEEtitleabstractindextext{%
\begin{abstract}
Predicting the eventual scale of information cascades during their early emergence is a fundamental challenge in online social network mining. In this paper, we present an end-to-end observational and predictive framework to model information diffusion across $960$ unique web media domains over $5,000$ inferred propagation edges from the MemeTracker NETINF cascade dataset. We formulate early cascade size prediction as a regression task over non-leaking temporal observation cutoffs ($10\%$, $20\%$, $30\%$, and $50\%$). By evaluating progressive feature additions, we demonstrate that combining early temporal dynamics with graph centrality metrics (Out-Degree, Betweenness, and PageRank) increases variance explained ($R^2$) from $0.2961$ to $0.3286$. Machine learning ensembles (Random Forest, Gradient Boosting, and XGBoost) achieve an $18.5\%$ reduction in Mean Absolute Error ($\text{MAE} = 8.55$ websites) compared to standard mean baselines. Furthermore, temporal link prediction heuristics achieve strong edge classification performance ($\text{ROC-AUC} = 0.7585$, $\text{PR-AUC} = 0.7549$). Our ablation study confirms that while temporal features serve as the primary baseline anchor, network centrality metrics provide the principal structural improvement.
\end{abstract}

\begin{IEEEkeywords}
Social Network Analysis, Information Cascades, Early Cascade Size Prediction, Community Detection, Network Centrality, Link Prediction, Temporal Leakage Audit.
\end{IEEEkeywords}}

\maketitle
\IEEEdisplaynontitleabstractindextext
\IEEEpeerreviewmaketitle

\section{Introduction}
\IEEEPARstart{I}{nformation} propagation across digital news platforms and social networks shapes public opinion, viral dynamics, and content dissemination. Understanding how a phrase, meme, or news story spreads—and predicting its eventual final reach—is vital for social media analytics, trend forecasting, and misinformation mitigation.

However, predicting cascade growth from early-stage behavior is notoriously difficult due to stochastic adoption times, heavy-tailed cascade distributions, and complex network topologies. A core research question is: \textit{Can network centrality, community structure, link-prediction features, and early temporal dynamics be combined to accurately predict the eventual final size of an information cascade?}

In this work, we implement a rigorous, reproducible Social Network Analysis (SNA) and machine learning pipeline on the MemeTracker NETINF dataset. We enforce strict scientific distinctions between observed temporal events, inferred propagation edges, and algorithmic community assignments.

\section{Related Work}
Diffusion network inference was pioneered by Gomez-Rodriguez et al.~\cite{gomez2010inferring}, who developed the NETINF algorithm to infer underlying hidden networks from timestamped meme cascades. Leskovec et al.~\cite{leskovec2009meme} analyzed phrase dynamics across news blogs, establishing heavy-tailed cascade distributions. Topological link prediction foundations were formalized by Liben-Nowell and Kleinberg~\cite{liben2007link}, while Blondel et al.~\cite{blondel2008fast} introduced the Louvain algorithm for fast community modularity optimization. Centrality metrics in social graphs originate from Freeman~\cite{freeman1978centrality} and Page et al.~\cite{page1999pagerank}. Ensemble machine learning methods like Random Forests~\cite{breiman2001random}, Gradient Boosting~\cite{friedman2001greedy}, and XGBoost~\cite{chen2016xgboost} have demonstrated strong performance in tabular network regression.

\section{Dataset \& Preprocessing Audit}
The project utilizes the MemeTracker NETINF inferred information diffusion network dataset (\texttt{InfoNet5000Q1000NEXP.txt}). The dataset comprises $5,000$ directed edge records inferring information flow across $Q=1000$ evaluated cascade trees.

\begin{table}[h]
\caption{Empirical Dataset & Preprocessing Audit Summary}
\label{tab:dataset_summary}
\centering
\begin{tabular}{ll}
\toprule
\textbf{Metric / Feature} & \textbf{Empirical Value} \\
\midrule
Raw Data Records & $5,000$ directed edges \\
Cleaned Data Records & $5,000$ (0 missing, 0 duplicates) \\
Unique Source Websites ($src$) & $733$ \\
Unique Destination Websites ($dst$) & $953$ \\
Total Unique Websites ($|V|$) & $960$ \\
Evaluated Cascade Trees & $1,000$ cascades \\
\bottomrule
\end{tabular}
\end{table}

\section{Network Construction \& Basic SNA}
We construct a Directed Weighted Graph $G = (V, E, W)$ where nodes represent website domain hostnames ($960$ nodes) and directed edges represent inferred propagation flows ($5,000$ edges). Edge weights encode cascade frequency (\texttt{number\_trees}), submodular marginal gain, and propagation delays.

\begin{table}[h]
\caption{Fundamental Social Network Analysis (SNA) Metrics}
\label{tab:network_metrics}
\centering
\begin{tabular}{ll}
\toprule
\textbf{SNA Metric} & \textbf{Empirical Result} \\
\midrule
Nodes ($|V|$) & $960$ \\
Edges ($|E|$) & $5,000$ \\
Network Density & $0.005431$ \\
Mean Total Degree & $10.4167$ \\
Max Out-Degree & $94$ (\texttt{startribune.com}) \\
Max In-Degree & $22$ (\texttt{sfgate.com}, \texttt{sott.net}) \\
Weakly Connected Components & $1$ ($100\%$ coverage) \\
Strongly Connected Components & $246$ (Largest SCC: $715$ nodes, $74.48\%$) \\
Directed Reciprocity & $0.0448$ ($4.48\%$) \\
Clustering Coefficient (Transitivity) & $0.020376$ ($2.04\%$) \\
Degree Assortativity & $-0.08954$ (Disassortative) \\
\bottomrule
\end{tabular}
\end{table}

\section{Centrality & Community Analysis}
We compute degree centrality, exact betweenness centrality, closeness centrality, and PageRank ($\alpha=0.85$). Top source hubs by out-degree include \texttt{startribune.com} ($0.0980$) and \texttt{smh.com.au} ($0.0761$), whereas top global PageRank influencers include \texttt{blogs.myspace.com} ($0.00581$) and \texttt{sott.net} ($0.00562$).

Louvain community detection partitions the $960$ nodes into $13$ detected network communities with a high modularity score of $Q = 0.4264$. Intra-community edges account for $43.44\%$ ($2,172$ edges), while inter-community edges account for $56.56\%$ ($2,828$ edges).

\section{Temporal Link Prediction Experiment}
We split the network chronologically into a historical training graph $G_{train}$ ($70\%$, $3,500$ edges) and future test edges ($30\%$, $1,500$ edges), balanced with $1,500$ sampled negative non-edges.

\begin{table}[h]
\caption{Link Prediction Classification Results}
\label{tab:link_prediction}
\centering
\begin{tabular}{lccccc}
\toprule
\textbf{Model} & \textbf{Precision} & \textbf{Recall} & \textbf{F1} & \textbf{ROC-AUC} & \textbf{PR-AUC} \\
\midrule
Logistic Regression & 0.7219 & 0.5711 & 0.6377 & 0.7455 & 0.7496 \\
Random Forest & 0.6921 & 0.7044 & 0.6982 & \textbf{0.7585} & \textbf{0.7549} \\
\bottomrule
\end{tabular}
\end{table}

\section{Primary Cascade Size Prediction}
The primary target variable is \texttt{target\_final_size} (total unique websites reached, mean = $623.71$). We partition $1,000$ cascades into $70\%$ train ($700$), $15\%$ validation ($150$), and $15\%$ test ($150$) at the cascade level to eliminate observation leakage.

\begin{table}[h]
\caption{Primary Machine Learning Model Evaluations (Test Set, $N=600$)}
\label{tab:ml_evaluations}
\centering
\begin{tabular}{lccc}
\toprule
\textbf{Model / Baseline} & \textbf{MAE (Websites)} & \textbf{RMSE} & \textbf{$R^2$ Score} \\
\midrule
Baseline: Mean Prediction & 10.4867 & 13.3194 & -0.0059 \\
Baseline: Median Prediction & 10.4867 & 13.3194 & -0.0059 \\
Baseline: Naive Early Size & 350.7750 & 368.4233 & -768.6243 \\
Ridge Regression & 9.8525 & 12.6087 & 0.0986 \\
XGBoost Regressor & 8.6229 & 11.0024 & 0.3136 \\
Gradient Boosting Regressor & \textbf{8.5516} & 10.9131 & 0.3247 \\
Random Forest Regressor & 8.6318 & \textbf{10.8724} & \textbf{0.3298} \\
\bottomrule
\end{tabular}
\end{table}

\section{Progressive Feature & Observation Dynamics}
Table~\ref{tab:progressive} details progressive feature additions. Adding centrality metrics to temporal features increases $R^2$ from $0.2961$ to $0.3286$.

\begin{table}[h]
\caption{Progressive Feature Addition Experiment}
\label{tab:progressive}
\centering
\begin{tabular}{lccc}
\toprule
\textbf{Feature Configuration} & \textbf{MAE} & \textbf{RMSE} & \textbf{$R^2$} \\
\midrule
Model A: Temporal Only (5 feat.) & 8.1634 & 10.4106 & 0.2961 \\
Model B: Model A + Centrality (11 feat.) & 7.9800 & \textbf{10.1673} & \textbf{0.3286} \\
Model C: Model B + Community (12 feat.) & 7.9883 & 10.2240 & 0.3211 \\
Model D: Full Model (15 feat.) & \textbf{7.9720} & 10.2041 & 0.3238 \\
\bottomrule
\end{tabular}
\end{table}

Evaluating observation depth ($p$) shows strong predictive scaling: at $10\%$ cutoff, $\text{MAE} = 8.57$, $R^2 = 0.2207$; at $50\%$ cutoff, $\text{MAE} = 7.06$, $R^2 = \mathbf{0.4753}$.

\section{Ablation & Feature Importance}
Ablation confirms temporal features anchor predictions ($R^2$ drops to $0.0319$ without them), while centrality provides structural gains ($R^2$ drops to $0.2996$ without centrality). Permutation feature importance ranks \texttt{early\_event\_count} ($0.7643$), \texttt{early\_size} ($0.2176$), and \texttt{observation\_level} ($0.1956$) as top predictors.

\section{Limitations & Conclusion}
Key limitations include observational edge inference (NETINF) rather than true causal tracking and compact final cascade size variance ($\sigma=13.08$). In conclusion, combining early temporal dynamics with network centrality metrics enables early cascade size prediction with high accuracy ($\text{MAE} = 8.55$ sites, $R^2 = 0.4753$ at $50\%$ observation).

\bibliographystyle{IEEEtran}
\bibliography{references}

\end{document}
